\begin{center}\begin{tikzpicture}[font=\small,>=stealth,line join=round]
\draw[gray!40,rounded corners=3pt] (0,0) rectangle (14.2,8.4);\node at (7.1,7.800000000000001) {比较的是直路 PA 与绕路 PBA，先得到范围界};
\begin{scope}[shift={(3.0,2.2)},x=1.2cm,y=1.2cm]\draw[->,gray] (-1.6,0)--(2.6,0) node[right]{$x$};\draw[->,gray] (0,-.35)--(0,2.8) node[left]{$y$};\draw[orange!85!black,line width=1.6pt] (0,2)--(-1,0);\draw[violet,thick] (2,1)--(-1,0);\draw[teal!80!black,thick] (0,2)--(2,1);\fill[teal!80!black] (0,2) circle(2.2pt) node[above left=4pt]{$A(0,2)$};\fill[teal!80!black] (2,1) circle(2.2pt) node[above right=4pt]{$B(2,1)$};\fill[blue] (-1,0) circle(2.2pt) node[below=4pt]{$P(-1,0)$};\end{scope}\node[anchor=west,align=left] at (8.6,4.4) {直路：$PA$\\[12pt]绕行：$PB+BA$\\[12pt]$PA\le PB+BA$\\[12pt]$PA-PB\le AB$\\[12pt]交换 $A,B$ 再得另一侧界。};
\node[text width=12.8cm,align=center] at (7.1,.65) {橙色为直路；紫色与青色合起来是绕路。这里只得到界，还未证明能取到。};
\end{tikzpicture}\end{center}
